\documentclass[12pt]{article}

\usepackage[a4paper, margin=2.5cm]{geometry}
\usepackage{amsmath}
\usepackage{amssymb}
\usepackage{amsthm}
\usepackage[colorlinks=true, linkcolor=blue, citecolor=blue, urlcolor=blue]{hyperref}
\usepackage{booktabs}
\usepackage{natbib}

\newtheorem{theorem}{Theorem}
\newtheorem{proposition}[theorem]{Proposition}
\newtheorem{corollary}[theorem]{Corollary}
\newtheorem{definition}[theorem]{Definition}
\newtheorem{remark}[theorem]{Remark}
\newtheorem{example}[theorem]{Example}

\newcommand{\Mcl}{\mathcal{M}_{\mathrm{cl}}}
\newcommand{\Fcal}{\mathcal{F}}
\newcommand{\Hcal}{\mathcal{H}}
\newcommand{\Lcal}{\mathcal{L}}
\newcommand{\Sadm}{\mathcal{S}_{\mathrm{adm}}}
\newcommand{\VW}{V_{\mathrm{W}}}

\title{Nelson Pairing for Closure Dynamics:\\Schr\"odinger-Type Evolution\\from Forward and Backward Closure Processes}
\author{%
  Lee Smart\\[4pt]
  \textit{Institute of Vibrational Field Dynamics}\\[2pt]
  \texttt{contact@vibrationalfielddynamics.org}\\[2pt]
  \texttt{@vfd\_org}%
}
\date{April 2026}

\begin{document}

\maketitle

\noindent\textbf{Status:} Preprint (not peer-reviewed)

\begin{abstract}
Paper~XVII proved that generators of the form $L_{\mathrm{BK}} + iM$ --- the backward Kolmogorov generator plus any imaginary potential --- are necessarily dissipative and cannot produce unitary evolution. It identified the Nelson-type forward/backward pairing as the most promising route past this obstruction. The present paper carries out that construction within the closure framework.

Starting from the forward closure Langevin process and its time-reversal, the current and osmotic velocities are defined and shown to produce an exact continuity equation. Because the closure drift is a gradient, the current velocity is automatically irrotational, and a phase function $S$ is determined without additional assumptions. The closure Madelung pair is derived: the continuity equation plus a modified Hamilton--Jacobi equation with an explicit quantum-potential-like term.

Recombining $\rho$ and $S$ into a Nelson wavefunction $\Psi_N = \sqrt{\rho}\,e^{iS/\sigma^2}$, the paper derives the \textit{closure-Nelson wave equation}: a nonlinear Schr\"odinger equation in which the potential depends on the state through $\Delta|\Psi_N|/|\Psi_N|$. The central result is two-tiered: evaluated on the Paper~XIV equilibrium solution, the equation becomes the Schr\"odinger equation with Witten Hamiltonian (and trivially satisfies $H_{\mathrm{W}}\Psi_{\mathrm{st}} = 0$); for perturbations around equilibrium, the equation formally reduces at leading order in a small-perturbation parameter $\varepsilon$ to the same linear Schr\"odinger dynamics, with remainder terms of order $O(\varepsilon)$ not rigorously estimated here.

This constitutes the first paired-process derivation of Schr\"odinger-type dynamics in the equilibrium and near-equilibrium regime within the closure framework: the imaginary kinetic term, absent in the single-generator framework of Papers~XV--XVII, emerges from the pairing of forward and backward processes. The remaining open problem is no longer the kinetic sector but the nonlinear state-dependent potential away from equilibrium. The resulting nonlinear correction is not interpreted here as a modification of quantum mechanics, but as a structurally derived residual whose physical status is deferred to subsequent work.
\end{abstract}

%==================================================================
\section{Introduction}\label{sec:intro}
%==================================================================

Paper~XVII~\citep{SmartXVII} proved the dissipative obstruction: any generator of the form $L_{\mathrm{BK}} + iM$, where $L_{\mathrm{BK}}$ is the backward Kolmogorov generator from the closure Langevin process, is dissipative in the equilibrium-weighted inner product regardless of the imaginary potential $iM$. This ruled out the single-generator approach of Papers~XV--XVI~\citep{SmartXV, SmartXVI} as a route to unitary evolution.

Paper~XVII classified three candidate routes past the obstruction and identified Route~B --- the Nelson-type forward/backward pairing --- as the most promising, because it uses the existing stochastic closure process and its time-reversal without introducing new canonical variables.

The present paper carries out this construction.

\subsection*{The problem}

The goal is not to assume Schr\"odinger structure, but to determine whether it \textit{emerges} from the paired forward/backward closure processes. Specifically:

\begin{enumerate}
    \item[(P10)] \textbf{The current-velocity problem.} Can the closure framework supply a non-stationary dynamical equation for the current velocity $v$?
    \item[(P11)] \textbf{The Schr\"odinger-recovery problem.} Does the Nelson wavefunction $\Psi_N = \sqrt{\rho}\,e^{iS/\sigma^2}$ satisfy a Schr\"odinger-type equation?
\end{enumerate}

\subsection*{Scope and epistemic status}

This paper derives the closure-Nelson wave equation and identifies the regime in which it reduces to the Schr\"odinger equation. Where results are derived, they are stated as propositions; where they depend on approximations (linearisation near equilibrium), that status is explicit. More precisely, the paper derives an exact nonlinear closure-Nelson wave equation and identifies the regime in which its leading-order dynamics reduces to the linear Schr\"odinger equation.

%==================================================================
\section{Forward and Backward Closure Processes}\label{sec:processes}
%==================================================================

\subsection{The forward process}

The closure Langevin process of Paper~XIV~\citep{SmartXIV} is:
\begin{equation}\label{eq:forward}
    d\psi_t = b_+(\psi_t)\, dt + \sigma\, dW_t^+, \qquad b_+ = -\nabla\Fcal
\end{equation}
where $\Fcal \in C^2(\Sadm)$ is the closure functional and $\sigma > 0$ is the noise amplitude. The probability density $\rho(\psi, t)$ satisfies the forward Fokker--Planck equation:
\begin{equation}\label{eq:FP}
    \partial_t \rho = \nabla\cdot(\rho\,\nabla\Fcal) + \frac{\sigma^2}{2}\Delta\rho
\end{equation}

The stationary distribution is $\rho_{\mathrm{st}} = Z^{-1}\exp(-2\Fcal/\sigma^2)$.

\subsection{The backward process}

The time-reversed diffusion has backward drift:
\begin{equation}\label{eq:backward_drift}
    b_-(\psi, t) = b_+(\psi) - \sigma^2\,\nabla\log\rho(\psi,t) = -\nabla\Fcal - \sigma^2\,\nabla\log\rho
\end{equation}

\subsection{Nelson velocities}

Define:
\begin{align}
    v &= \frac{b_+ + b_-}{2} = -\nabla\Fcal - \frac{\sigma^2}{2}\nabla\log\rho \qquad \text{(current velocity)} \label{eq:v} \\
    u &= \frac{b_+ - b_-}{2} = \frac{\sigma^2}{2}\nabla\log\rho \qquad \text{(osmotic velocity)} \label{eq:u}
\end{align}

The current velocity represents the net probabilistic flow; the osmotic velocity represents the diffusive spreading.

\begin{remark}[Stationarity check]\label{rem:stationary}
At the Paper~XIV equilibrium $\rho = \rho_{\mathrm{st}}$: $\nabla\log\rho_{\mathrm{st}} = -2\nabla\Fcal/\sigma^2$, giving $v_{\mathrm{st}} = 0$ and $u_{\mathrm{st}} = -\nabla\Fcal$. The vanishing current velocity at equilibrium is consistent with the ground-state condition in Nelson mechanics.
\end{remark}

%==================================================================
\section{Continuity Equation}\label{sec:continuity}
%==================================================================

\begin{proposition}[Continuity equation]\label{prop:continuity}
The probability density $\rho$ and current velocity $v$ satisfy:
\begin{equation}\label{eq:continuity}
    \partial_t\rho + \nabla\cdot(\rho\, v) = 0
\end{equation}
\end{proposition}

\begin{proof}
From~\eqref{eq:v}, $\rho\, v = -\rho\,\nabla\Fcal - (\sigma^2/2)\nabla\rho$, so:
\begin{equation}
    \nabla\cdot(\rho\, v) = -\nabla\cdot(\rho\,\nabla\Fcal) - \frac{\sigma^2}{2}\Delta\rho
\end{equation}
Comparing with the Fokker--Planck equation~\eqref{eq:FP}: $\partial_t\rho = \nabla\cdot(\rho\,\nabla\Fcal) + (\sigma^2/2)\Delta\rho = -\nabla\cdot(\rho\,v)$.
\end{proof}

This is the first Madelung equation. It is exact and requires no approximation. It is also the first place where the paired-process construction departs structurally from the single-generator framework of Paper~XVI: the current velocity $v$ involves both forward and backward drifts, not just the forward drift $b_+$.

%==================================================================
\section{Phase Function and Irrotationality}\label{sec:phase}
%==================================================================

\begin{proposition}[Irrotational current velocity]\label{prop:irrotational}
On the open subset $\{(\psi, t) : \rho(\psi, t) > 0\}$ of configuration space--time (i.e., away from nodal sets of the density), the current velocity $v$ defined in~\eqref{eq:v} is a gradient field:
\begin{equation}\label{eq:v_gradient}
    v = \nabla S, \qquad S(\psi,t) = -\Fcal(\psi) - \frac{\sigma^2}{2}\log\rho(\psi,t).
\end{equation}
\end{proposition}

\begin{proof}
On the positivity region $\{\rho > 0\}$, $\log\rho$ is defined and differentiable, so $-(\sigma^2/2)\nabla\log\rho$ is a gradient; $-\nabla\Fcal$ is a gradient wherever $\Fcal \in C^1$. Their sum is therefore a gradient. The phase function $S$ is determined (up to a spatially constant, time-dependent function) by $\nabla S = v$. The restriction to $\rho > 0$ is load-bearing: nodal sets of $\rho$ are genuine obstructions to global gradient-field representation.
\end{proof}

\begin{remark}
Irrotationality is not an assumption --- it is a consequence of the closure framework, in which both the drift $b_+ = -\nabla\Fcal$ and the density correction $\nabla\log\rho$ are gradient fields. Because irrotationality is inherited from the gradient structure of both the closure drift and the density correction, no separate phase ansatz is needed: the phase function $S$ is forced by the closure geometry. This is a structural advantage of the closure framework over generic stochastic processes, where $v$ need not be irrotational.
\end{remark}

Writing $R = \sqrt{\rho}$, the phase can be expressed as:
\begin{equation}\label{eq:S}
    S = -\Fcal - \sigma^2\log R
\end{equation}

%==================================================================
\section{The Closure Madelung Equations}\label{sec:madelung}
%==================================================================

We now derive the second Madelung equation from the Fokker--Planck dynamics.

\begin{proposition}[Closure Madelung pair]\label{prop:madelung}
The pair $(\rho, S)$ satisfies:
\begin{align}
    \partial_t\rho + \nabla\cdot(\rho\,\nabla S) &= 0 \label{eq:M1} \\[4pt]
    \partial_t S + \frac{|\nabla S|^2}{2} - \sigma^2 \VW + \frac{\sigma^4}{2}\frac{\Delta\sqrt{\rho}}{\sqrt{\rho}} &= 0 \label{eq:M2}
\end{align}
where $\VW = |\nabla\Fcal|^2/(2\sigma^2) - \Delta\Fcal/2$ is the Witten potential of Paper~XVI.
\end{proposition}

\begin{proof}
Equation~\eqref{eq:M1} is Proposition~\ref{prop:continuity} with $v = \nabla S$.

For~\eqref{eq:M2}: from $S = -\Fcal - \sigma^2\log R$ with $R = \sqrt{\rho}$, the time derivative is $\partial_t S = -\sigma^2\partial_t\log R = -(\sigma^2/2)\partial_t\log\rho$. From the Fokker--Planck:
\begin{equation}
    \frac{\partial_t\rho}{\rho} = \Delta\Fcal + \nabla\Fcal\cdot\nabla\log\rho + \frac{\sigma^2}{2}(\Delta\log\rho + |\nabla\log\rho|^2)
\end{equation}

Using $\nabla\log\rho = 2\nabla\log R$ and $\Delta\log\rho = 2\Delta\log R$:
\begin{equation}
    \partial_t S = -\frac{\sigma^2}{2}\Delta\Fcal - \sigma^2\nabla\Fcal\cdot\nabla\log R - \frac{\sigma^4}{2}\Delta\log R - \sigma^4|\nabla\log R|^2
\end{equation}

The kinetic term is $|v|^2/2 = |\nabla\Fcal|^2/2 + \sigma^2\nabla\Fcal\cdot\nabla\log R + (\sigma^4/2)|\nabla\log R|^2$.

Adding, and using $\Delta\log R + |\nabla\log R|^2 = \Delta R/R$:
\begin{align}
    \partial_t S + \frac{|v|^2}{2} &= \frac{|\nabla\Fcal|^2}{2} - \frac{\sigma^2}{2}\Delta\Fcal - \frac{\sigma^4}{2}\frac{\Delta R}{R} \notag \\
    &= \sigma^2 \VW - \frac{\sigma^4}{2}\frac{\Delta\sqrt{\rho}}{\sqrt{\rho}}
\end{align}
Rearranging gives~\eqref{eq:M2}.
\end{proof}

\subsection{Comparison with the Schr\"odinger Madelung equations}

The standard Schr\"odinger equation $i\hbar\partial_t\psi = -(\hbar^2/2m)\Delta\psi + V\psi$ is equivalent to the Madelung pair:
\begin{align}
    \partial_t\rho + \nabla\cdot(\rho\,\nabla S/m) &= 0 \label{eq:SM1}\\
    \partial_t S + \frac{|\nabla S|^2}{2m} + V - \frac{\hbar^2}{2m}\frac{\Delta\sqrt\rho}{\sqrt\rho} &= 0 \label{eq:SM2}
\end{align}

The closure Madelung equation~\eqref{eq:M2} has the same structure as~\eqref{eq:SM2} with:
\begin{itemize}
    \item effective mass $m_{\mathrm{eff}} = 1$ (from $|v|^2/2$ rather than $|v|^2/(2m)$),
    \item effective potential $V_{\mathrm{eff}} = -\sigma^2\VW$ (with opposite sign to the usual convention),
    \item quantum potential $Q = +(\sigma^4/2)\Delta\sqrt\rho/\sqrt\rho$ (with \textbf{opposite sign} to the Schr\"odinger convention).
\end{itemize}

The sign reversal of the quantum potential is the key structural difference between the closure Madelung system and the standard Schr\"odinger Madelung system. However, this should not be read as a failure of the pairing construction. Rather, it identifies the precise way in which the closure-paired dynamics differs from exact linear quantum mechanics away from equilibrium: the sign discrepancy is absorbed into a state-dependent potential in the recombined wave equation (Section~\ref{sec:wave}), drops out exactly at equilibrium, and is controlled perturbatively --- at leading order in the small parameter $\varepsilon$ --- in the near-equilibrium linearisation (Section~\ref{sec:linear}).

%==================================================================
\section{The Closure-Nelson Wave Equation}\label{sec:wave}
%==================================================================

Recombining $\rho$ and $S$ into the Nelson wavefunction:
\begin{equation}\label{eq:nelson_psi}
    \Psi_N = \sqrt{\rho}\, e^{iS/\sigma^2}
\end{equation}

\begin{proposition}[Closure-Nelson wave equation]\label{prop:wave}
The Madelung pair~\eqref{eq:M1}--\eqref{eq:M2} is equivalent to the nonlinear wave equation:
\begin{equation}\label{eq:closure_schrodinger}
    \boxed{i\sigma^2\,\partial_t\Psi_N = -\frac{\sigma^4}{2}\Delta\Psi_N + U[\Psi_N]\,\Psi_N}
\end{equation}
where the state-dependent potential is:
\begin{equation}\label{eq:U}
    U[\Psi_N] = -\sigma^2\VW + \sigma^4\frac{\Delta|\Psi_N|}{|\Psi_N|}
\end{equation}
\end{proposition}

\begin{proof}
The standard Madelung correspondence states that $i\sigma^2\partial_t\Psi = -(\sigma^4/2)\Delta\Psi + U\Psi$ with $\Psi = Re^{iS/\sigma^2}$ is equivalent to:
\begin{align}
    &\partial_t\rho + \nabla\cdot(\rho\,\nabla S) = 0 \\
    &\partial_t S + \frac{|\nabla S|^2}{2} + U - \frac{\sigma^4}{2}\frac{\Delta R}{R} = 0
\end{align}
Comparing the second equation with~\eqref{eq:M2}:
\begin{equation}
    U - \frac{\sigma^4}{2}\frac{\Delta R}{R} = -\sigma^2\VW + \frac{\sigma^4}{2}\frac{\Delta R}{R}
\end{equation}
Solving: $U = -\sigma^2\VW + \sigma^4\Delta R/R = -\sigma^2\VW + \sigma^4\Delta|\Psi_N|/|\Psi_N|$.
\end{proof}

\begin{remark}[Nature of the equation]
Equation~\eqref{eq:closure_schrodinger} is a \textit{nonlinear} Schr\"odinger equation: the potential $U$ depends on the state $\Psi_N$ through $\Delta|\Psi_N|/|\Psi_N|$. It has the correct kinetic structure ($-(\sigma^4/2)\Delta$ with the imaginary prefactor $i\sigma^2$) but a state-dependent potential. The nonlinearity enters through the quantum-potential sign reversal identified in Section~\ref{sec:madelung}.

The paired-process construction produces an imaginary kinetic term $-(\sigma^4/2)\Delta$ in the recombined wave equation, which was impossible in the single-generator setting of Papers~XVI--XVII. What remains is not the kinetic problem but the residual state dependence of the effective potential.
\end{remark}

\begin{remark}[Why this evades the XVII obstruction]\label{rem:obstruction}
Paper~XVII ruled out unitarity for generators of the single-process Kolmogorov form $L_{\mathrm{BK}} + iM$. The present construction does not modify that form; it leaves it behind. The forward/backward pairing produces a different composite dynamics: the current velocity $v = (b_+ + b_-)/2$ involves both forward and backward drifts, and the resulting kinetic structure is no longer governed by a single dissipative Kolmogorov generator. The imaginary kinetic term $-(\sigma^4/2)\Delta$ therefore emerges from the pairing structure itself, not from any modification of a single Kolmogorov generator, and is consistent with the obstruction theorem of Paper~XVII.
\end{remark}

%==================================================================
\section{Geometric Structure of the Nelson Wavefunction}\label{sec:geometric}
%==================================================================

\begin{remark}[Addendum status]\label{rem:addendum_geometric}
This section was added in a subsequent revision to make explicit a bookkeeping observation about the Nelson-pairing construction: the multiplication-by-$i$ that takes amplitude $\sqrt\rho$ to phase-weighted tangents $(R/\sigma^2)\delta S$ is the standard complex structure of the Hilbert space $L^2$, read through the amplitude--phase chart constructed in the rest of this paper. The wavefunction $\Psi_N = \sqrt{\rho}\,e^{iS/\sigma^2}$ of~\eqref{eq:nelson_psi} and the wave equation of Proposition~\ref{prop:wave} are unchanged. The observation is cited by Paper~XXI (which combines it with a symplectic form to obtain a K\"ahler structure at the equilibrium tangent) and by Paper~XXX (which uses the K\"ahler structure in its Born-rule derivation). No independent geometric content beyond the standard complex structure of $L^2$ is claimed here.
\end{remark}

The Nelson wavefunction~\eqref{eq:nelson_psi} recombines two real-valued fields --- the amplitude density $\sqrt\rho$ and the phase function $S$ --- into a single complex-valued object. This section records, in amplitude--phase coordinates, how multiplication by $i$ acts on tangent vectors to admissible wavefunctions, and identifies the regime (equilibrium tangent) in which this action coincides with the standard complex structure of a linear complex vector space.

\subsection{The amplitude-phase pairing}\label{subsec:amp_phase}

Every admissible Nelson wavefunction decomposes, at each point $\psi$ in the base space of configurations where $\rho(\psi, t) > 0$, as
\begin{equation}\label{eq:amp_phase_pair}
    \Psi_N(\psi) \;=\; R(\psi)\, e^{iS(\psi)/\sigma^2},
    \qquad R(\psi) = \sqrt{\rho(\psi)} \in \mathbb{R}_{>0},
    \qquad S(\psi) \in \mathbb{R}.
\end{equation}
At a fixed reference wavefunction $\Psi_{N,0}$ (for example, the Paper~XIV stationary state $\Psi_{\mathrm{st}}$), a tangent vector $\delta\Psi_N$ splits into an amplitude variation $\delta R$ and a phase variation $\delta S$:
\begin{equation}\label{eq:tangent_split}
    \delta\Psi_N(\psi) \;=\; \Bigl(\delta R(\psi) \;+\; \frac{i R_0(\psi)}{\sigma^2}\, \delta S(\psi)\Bigr) e^{iS_0(\psi)/\sigma^2}.
\end{equation}
Equation~\eqref{eq:tangent_split} is the infinitesimal form of the Nelson pairing: tangent vectors carry one real (amplitude) and one imaginary (phase) component, weighted by $R_0(\psi)/\sigma^2$. The decomposition is defined on the region $\{\psi : \rho_0(\psi) > 0\}$ of the base space; nodal points of the reference state are excluded from the present analysis.

\subsection{Multiplication-by-$i$ as an endomorphism of tangent vectors}\label{subsec:J}

Define the endomorphism $J$ on tangent vectors by multiplication by $i$ on the right-hand side of~\eqref{eq:tangent_split}:
\begin{equation}\label{eq:J_def}
    J\,\delta\Psi_N \;:=\; i\, \delta\Psi_N
    \;=\; \Bigl(-\frac{R_0(\psi)}{\sigma^2}\,\delta S(\psi) \;+\; i\,\delta R(\psi)\Bigr) e^{iS_0(\psi)/\sigma^2}.
\end{equation}
In amplitude--phase coordinates, $J$ acts as
\begin{equation}\label{eq:J_matrix}
    J: (\delta R, \delta S) \;\longmapsto\; \Bigl(-\frac{R_0}{\sigma^2}\,\delta S,\; \frac{\sigma^2}{R_0}\,\delta R\Bigr),
\end{equation}
so that $J^2 = -\mathrm{id}$. Squaring reverses sign, and $J$ rotates amplitude variations into phase variations.

\begin{remark}[Observation, not a theorem]\label{prop:J}
Equations~\eqref{eq:J_def}--\eqref{eq:J_matrix} record a bookkeeping observation: in the Nelson amplitude--phase parameterisation, multiplication by $i$ on the complex wavefunction $\Psi_N = R\,e^{iS/\sigma^2}$ is nothing other than rotating amplitude variations into phase variations (up to the scale factor $R_0/\sigma^2$ set by the reference state). $J$ here is the standard complex structure of the linear complex vector space in which $\Psi_N$ takes values, read through the amplitude--phase chart. We do not claim a new theorem: $J^2 = -\mathrm{id}$ is tautological once a complex-valued wavefunction has been defined.

What the Nelson pairing contributes, beyond a complex wavefunction, is the \emph{interpretation} of the real and imaginary parts as amplitude and phase of the closure-stochastic density, with both parts determined by the forward/backward process construction of Section~\ref{sec:processes}. The pairing therefore explains \emph{why} amplitude and phase are to be paired as real and imaginary components of the same object --- a question not addressed by simply writing down a complex wavefunction --- but does not add new geometric content beyond the standard complex structure of $L^2$.
\end{remark}

\subsection{Regime in which the linear complex structure applies}\label{subsec:J_integrable}

\begin{remark}[Linear complex structure at the equilibrium tangent]\label{rem:J_integrable}
At the equilibrium state $\Psi_{\mathrm{st}}$, the linearised wave equation of Proposition~\ref{prop:linearisation} is $i\,\partial_t\Psi = H_{\mathrm{W}}\Psi$, a linear Schr\"odinger-type evolution on a linear complex vector space. On this linear space, the standard complex structure (multiplication by $i$) is integrable by construction --- the space is globally a complex vector space, not just pointwise --- and $J$ as defined in~\eqref{eq:J_def} coincides with this standard complex structure on the amplitude--phase coordinates of the tangent at the equilibrium point. No nontrivial integrability claim is made; the content is that equilibrium-tangent perturbations of $\Psi_{N,0}$ live naturally in a linear complex vector space with $J$ as its complex structure.
\end{remark}

\begin{remark}[Beyond the equilibrium tangent]\label{rem:J_nonlinear}
Away from the equilibrium tangent, the nonlinear residual of Proposition~\ref{prop:wave} perturbs the dynamics. The pointwise decomposition~\eqref{eq:J_def} still makes sense on the nonvanishing-amplitude region, but defining a complex-manifold structure on the full admissible wavefunction space would require additional topological and differentiable setup that is not undertaken here. The present section therefore makes no claim about geometric structure beyond the equilibrium-tangent linearisation.
\end{remark}

\subsection{Role in the downstream programme}\label{subsec:xviii_role}

\begin{remark}[How Papers~XXI and XXX use this observation]\label{rem:J_downstream}
Paper~XXI takes the multiplication-by-$i$ structure at the equilibrium tangent as an input and combines it with the symplectic form induced by the linearised Schr\"odinger dynamics to exhibit a K\"ahler structure on the projective space of equilibrium-tangent states (see Paper~XXI for the construction and scope). Paper~XXX uses this K\"ahler structure, together with the closure Hessian providing a Riemannian normalisation, to identify the admissible-state geometry with Fubini--Study up to scale via Petz's uniqueness theorem. The present section supplies the amplitude--phase interpretation of multiplication by $i$; the deeper geometric content resides in Papers~XXI and~XXX.
\end{remark}

%==================================================================
\section{Near-Equilibrium Linearisation}\label{sec:linear}
%==================================================================

The nonlinear potential $U[\Psi_N]$ simplifies dramatically near the Paper~XIV equilibrium.

\begin{proposition}[Equilibrium and near-equilibrium Schr\"odinger regime]\label{prop:linearisation}
The closure-Nelson wave equation~\eqref{eq:closure_schrodinger} admits a Schr\"odinger regime in two stages:
\begin{enumerate}
    \item \textbf{Exact on the equilibrium solution.} At $\rho = \rho_{\mathrm{st}}$, the state-dependent potential $U[\Psi]$ reduces to the $\psi$-dependent function
    \begin{equation}\label{eq:U_stationary}
        U[\Psi_{\mathrm{st}}](\psi) = \sigma^2\VW(\psi),
    \end{equation}
    which (despite not being constant in $\psi$) is now state-independent. Evaluated on the equilibrium solution itself, the wave equation reduces to $i\sigma^2\partial_t\Psi_{\mathrm{st}} = \sigma^2 H_{\mathrm{W}}\Psi_{\mathrm{st}}$; since $H_{\mathrm{W}}\Psi_{\mathrm{st}} = 0$ for a normalised stationary ground state (Paper~XVI), this is consistent with the time-independence of $\Psi_{\mathrm{st}}$.

    \item \textbf{Leading-order near equilibrium.} For perturbations $\Psi = \Psi_{\mathrm{st}}(1 + \varepsilon\eta)$ with $\varepsilon \ll 1$, the nonlinear correction to the potential is $U[\Psi] = \sigma^2\VW + O(\varepsilon)$; the equation for the perturbation $\eta$ at leading order in $\varepsilon$ reduces formally to
    \begin{equation}\label{eq:linearised}
        \boxed{i\,\partial_t\Psi = H_{\mathrm{W}}\,\Psi}
    \end{equation}
    where $H_{\mathrm{W}} = -(\sigma^2/2)\Delta + \VW$ is the non-negative Witten Hamiltonian (non-negativity is proved in Paper~XVII, not XVI, via supersymmetric factorisation). The word ``linearises'' refers to the leading-order expansion; the remainder terms are formally $O(\varepsilon)$ but are not estimated with remainder bounds here.
\end{enumerate}
\end{proposition}

\begin{proof}
At stationarity, $R_{\mathrm{st}} = \sqrt{\rho_{\mathrm{st}}} \propto e^{-\Fcal/\sigma^2}$, so:
\begin{equation}
    \frac{\Delta R_{\mathrm{st}}}{R_{\mathrm{st}}} = \frac{|\nabla\Fcal|^2}{\sigma^4} - \frac{\Delta\Fcal}{\sigma^2} = \frac{2\VW}{\sigma^2}
\end{equation}
Substituting into~\eqref{eq:U}: $U[\Psi_{\mathrm{st}}] = -\sigma^2\VW + \sigma^4 \cdot 2\VW/\sigma^2 = \sigma^2\VW$.

The full wave equation~\eqref{eq:closure_schrodinger} at stationarity is:
\begin{equation}
    i\sigma^2\partial_t\Psi = -\frac{\sigma^4}{2}\Delta\Psi + \sigma^2\VW\,\Psi = \sigma^2\!\left(-\frac{\sigma^2}{2}\Delta + \VW\right)\!\Psi = \sigma^2 H_{\mathrm{W}}\Psi
\end{equation}
Dividing by $\sigma^2$ gives $i\partial_t\Psi = H_{\mathrm{W}}\Psi$.

For the linearisation: write $\Psi = \Psi_{\mathrm{st}}(1 + \varepsilon\eta)$ for small perturbation $\eta$. Then $|\Psi| = R_{\mathrm{st}}(1 + \varepsilon\,\mathrm{Re}(\eta)) + O(\varepsilon^2)$, and:
\begin{equation}
    \frac{\Delta|\Psi|}{|\Psi|} = \frac{\Delta R_{\mathrm{st}}}{R_{\mathrm{st}}} + O(\varepsilon) = \frac{2\VW}{\sigma^2} + O(\varepsilon)
\end{equation}
So $U = \sigma^2\VW + O(\varepsilon)$, and the linearised equation is $i\partial_t\Psi = H_{\mathrm{W}}\Psi$ to leading order.
\end{proof}

Thus the exact equilibrium state furnishes both the ground state of the closure dynamics and the base point around which a formal linear Schr\"odinger dynamics emerges at leading order (at the equilibrium tangent).

\begin{remark}[Interpretation]
The Witten Hamiltonian $H_{\mathrm{W}}$ appeared in Paper~XVI as the \textit{real part} of the drift-free generator, where it produced dissipation. Here, via the Nelson pairing, it reappears as the \textit{quantum Hamiltonian} governing oscillatory dynamics. The same operator that drives dissipation in the single-process framework drives oscillation in the paired-process framework.
\end{remark}

\begin{remark}[On the non-triviality of the linear limit]
It is true that any nonlinear dynamical system admits a linearisation at a fixed point. The non-trivial result here is not the existence of a linearisation, but the specific structure of that linearisation: the linearised generator is exactly the Witten Hamiltonian $H_{\mathrm{W}}$, which is derived from the closure functional $\Fcal$ and the noise scale $\sigma$ rather than chosen or postulated. The content is therefore not merely that a linear regime exists, but that it reproduces a specific quantum Hamiltonian determined by the closure geometry.
\end{remark}

\begin{remark}[Effective parameters]\label{rem:parameters}
Within this normalisation, the linearised equation $i\partial_t\Psi = H_{\mathrm{W}}\Psi$ may be written in Schr\"odinger form with effective parameters:
\begin{itemize}
    \item $\hbar_{\mathrm{eff}} = 1$ in these units (or $\sigma^2$ before the division),
    \item effective mass $m_{\mathrm{eff}} = 1/\sigma^2$ (from $-(\sigma^2/2)\Delta = -\hbar_{\mathrm{eff}}^2/(2m_{\mathrm{eff}})\Delta$),
    \item effective potential $\VW = |\nabla\Fcal|^2/(2\sigma^2) - \Delta\Fcal/2$.
\end{itemize}
These are read off from the closure functional $\Fcal$ and the noise scale $\sigma$; no additional postulates are needed within the linearised regime.
\end{remark}

%==================================================================
\section{Worked Example: Quadratic Closure Functional}\label{sec:example}
%==================================================================

\begin{example}[Quadratic closure]\label{ex:quadratic}
Let $\Fcal(\psi) = \frac{1}{2}k\psi^2$ in one dimension.

\subsection*{Nelson velocities}

$b_+ = -k\psi$, $\rho_{\mathrm{st}} = Z^{-1}e^{-k\psi^2/\sigma^2}$, $\nabla\log\rho_{\mathrm{st}} = -2k\psi/\sigma^2$.

At stationarity: $v = 0$, $u = -k\psi$. The phase $S = -k\psi^2/2 - \sigma^2\log R_{\mathrm{st}} = 0$ (constant, since $R_{\mathrm{st}} = Z^{-1/2}e^{-k\psi^2/(2\sigma^2)}$ and $-k\psi^2/2 - \sigma^2(-k\psi^2/(2\sigma^2)) = 0$).

\subsection*{Witten potential and Hamiltonian}

$\VW = k^2\psi^2/(2\sigma^2) - k/2$. The Witten Hamiltonian is:
\begin{equation}
    H_{\mathrm{W}} = -\frac{\sigma^2}{2}\partial_\psi^2 + \frac{k^2\psi^2}{2\sigma^2} - \frac{k}{2}
\end{equation}
This is a shifted quantum harmonic oscillator with spectrum $E_n = nk$, $n = 0, 1, 2, \ldots$

\subsection*{Linearised Schr\"odinger dynamics}

The linearised equation $i\partial_t\Psi = H_{\mathrm{W}}\Psi$ has eigenstates $\phi_n$ satisfying $H_{\mathrm{W}}\phi_n = nk\,\phi_n$. The ground state $\phi_0 \propto e^{-k\psi^2/(2\sigma^2)}$ is stationary. Excited states oscillate:
\begin{equation}
    \Psi_n(t) = e^{-inkt}\phi_n
\end{equation}
with fundamental frequency $\omega = k$. The energy gap is $\Delta E = k$ --- determined entirely by the closure functional.

\subsection*{Comparison with Paper~XVII}

In Paper~XVII's dissipative framework, the same excited states $\phi_n$ dissipate to the ground state at rate $2nk$. In the Nelson-paired framework, they oscillate with frequency $nk$. The Witten Hamiltonian governs both: dissipation rate $= 2E_n$, oscillation frequency $= E_n$.

\begin{table}[ht]
\centering
\caption{Quadratic example: dissipative vs.\ oscillatory dynamics.}
\label{tab:comparison}
\begin{tabular}{@{}lll@{}}
\toprule
Feature & Paper XVII (single process) & Paper XVIII (paired processes) \\
\midrule
Generator & $-H_{\mathrm{W}} + iM$ & $iH_{\mathrm{W}}$ (linearised) \\
Ground state & $e^{-k\psi^2/(2\sigma^2)}$, stable & $e^{-k\psi^2/(2\sigma^2)}$, stable \\
Excited states & Dissipate at rate $2nk$ & Oscillate at frequency $nk$ \\
Norm & Decreasing & Preserved (linearised) \\
\bottomrule
\end{tabular}
\end{table}

The quadratic example shows the programme transition in its cleanest form: the same Witten operator that appears as a dissipation generator in the single-process framework (Paper~XVII) becomes the oscillation generator in the paired-process linearised regime.
\end{example}

%==================================================================
\section{Conditions for Exact Schr\"odinger Recovery}\label{sec:conditions}
%==================================================================

The closure-Nelson wave equation~\eqref{eq:closure_schrodinger} is exactly Schr\"odinger (i.e.\ the potential $U$ is state-independent) if and only if:
\begin{equation}
    \sigma^4\frac{\Delta|\Psi_N|}{|\Psi_N|} = \text{function of $\psi$ only, independent of $t$}
\end{equation}

This is a strong condition. In the present framework, the physically relevant instance occurs in the near-equilibrium regime: for small perturbations $\Psi = \Psi_{\mathrm{st}}(1 + \varepsilon\eta)$, the nonlinear term is $O(\varepsilon)$ and the linearised dynamics is exactly Schr\"odinger (Proposition~\ref{prop:linearisation}). For generic non-equilibrium states, the nonlinear correction is present. The correction can be characterised as:
\begin{equation}
    \delta U = \sigma^4\left(\frac{\Delta|\Psi_N|}{|\Psi_N|} - \frac{\Delta R_{\mathrm{st}}}{R_{\mathrm{st}}}\right) = \sigma^4\left(\frac{\Delta R}{R} - \frac{2\VW}{\sigma^2}\right)
\end{equation}
This vanishes when $R$ has the same curvature structure as $R_{\mathrm{st}}$ --- i.e.\ when the density profile is close to the equilibrium Gaussian (in the quadratic case) or more generally when $\rho$ is close to $\rho_{\mathrm{st}}$. Accordingly, the generic paired closure dynamics is nonlinear, while exact linear Schr\"odinger evolution appears as a distinguished regime selected by equilibrium structure or special modulus profiles.

%==================================================================
\section{Interpretation of the Nonlinear Residual}\label{sec:residual}
%==================================================================

The closure-Nelson wave equation~\eqref{eq:closure_schrodinger} differs from the linear Schr\"odinger equation with Witten Hamiltonian $H_{\mathrm{W}}$ (the target of the equilibrium-tangent reduction, Proposition~\ref{prop:linearisation}) by the equilibrium-centred modulus-curvature term
\begin{equation}\label{eq:residual}
    \delta U[\Psi_N] \;=\; \sigma^4\!\left(\frac{\Delta|\Psi_N|}{|\Psi_N|} - \frac{\Delta R_{\mathrm{st}}}{R_{\mathrm{st}}}\right),
\end{equation}
defined as the deviation from the equilibrium modulus curvature (equation~\eqref{eq:residual} above agrees with the deviation form used at the end of Section~\ref{sec:linear}). This residual vanishes exactly at equilibrium, $\delta U[\Psi_{\mathrm{st}}] = 0$, and is $O(\varepsilon)$ for perturbations $\Psi = \Psi_{\mathrm{st}}(1 + \varepsilon\eta)$ with $\varepsilon$ small. The raw modulus-curvature functional $\sigma^4 \Delta|\Psi_N|/|\Psi_N|$ that appears unsubtracted in~\eqref{eq:U} is \emph{not} the residual: at equilibrium it takes the nonzero value $2\sigma^2 \VW$, which is exactly what the subtraction in~\eqref{eq:residual} removes to produce the equilibrium-centred $\delta U$.

\subsection*{Structural origin}

The residual term arises from the same quantity that produces the quantum potential in the Madelung formulation:
\begin{equation}
    Q = \frac{\hbar^2}{2m}\frac{\Delta\sqrt{\rho}}{\sqrt{\rho}}
\end{equation}

In standard quantum mechanics, $Q$ appears with a sign that cancels nonlinearity when recombined into the Schr\"odinger equation. In the closure framework, the sign inherited from the dissipative forward process leads to a non-cancelling contribution, producing the residual $\delta U$.

\subsection*{What the residual measures}

The residual encodes the deviation of the density profile from the equilibrium structure:
\begin{equation}
    \delta U = \sigma^4 \left(\frac{\Delta R}{R} - \frac{\Delta R_{\mathrm{st}}}{R_{\mathrm{st}}}\right)
\end{equation}

Thus:
\begin{itemize}
    \item $\delta U = 0$ for equilibrium density profiles,
    \item $\delta U \neq 0$ for non-equilibrium configurations,
    \item the magnitude of $\delta U$ encodes the curvature deviation of $\rho$ from $\rho_{\mathrm{st}}$.
\end{itemize}

\subsection*{Status}

At this stage, the residual term is identified as a structural consequence of the closure pairing, not interpreted as a modification of quantum mechanics. Whether it represents:
\begin{itemize}
    \item a physical nonlinear correction to linear quantum mechanics,
    \item a gauge artefact removable by transformation,
    \item or a structurally meaningful quantity with independent physical content,
\end{itemize}
remains an open question and is deferred to subsequent work. In this sense, the residual is best understood at present as the closure-sector remainder after linear quantum behaviour has been recovered, rather than as a defect of the paired construction.

%==================================================================
\section{Reconciliation with Papers~XV--XVII}\label{sec:reconciliation}
%==================================================================

\subsection{Paper~XV: phase and interference}

Paper~XV introduced the complexified path integral with amplitude $\mathcal{A} = e^{-S/\sigma^2 + i\Theta/\sigma^2}$. The Nelson wavefunction $\Psi_N = \sqrt\rho\, e^{iS/\sigma^2}$ is a different object: it is constructed from the density and phase of the paired processes, not from the path-integral amplitude. However, at equilibrium both reduce to $\Psi \propto e^{-\Fcal/\sigma^2}$ (the same ground state), and both produce the squared-modulus probability structure $P = |\Psi|^2 = \rho$.

\subsection{Paper~XVI: evolution equation}

The closure evolution equation of Paper~XVI, $\partial_t\Psi = ({\sigma^2}/{2})\Delta\Psi - \nabla\Fcal\cdot\nabla\Psi + (i/\sigma^2)\Fcal\Psi$, governs the single forward process. The Nelson wave equation~\eqref{eq:closure_schrodinger} governs the paired forward-backward system. They share the same equilibrium state but have different generator structures: the former is dissipative (Paper~XVII), the latter is oscillatory (in the linearised regime).

\subsection{Paper~XVII: dissipative obstruction}

The obstruction theorem applies to generators of the single-process form $L_{\mathrm{BK}} + iM$. The Nelson construction evades it by not using a single Kolmogorov generator: the current velocity $v = (b_+ + b_-)/2$ involves both forward and backward drifts, and the resulting dynamics is not of the $L_{\mathrm{BK}} + iM$ form. The imaginary kinetic term emerges from the pairing, not from a modification of the single-process generator. In this sense, XVIII should be read not as a contradiction of XVII, but as the constructive completion of the route XVII left open.

%==================================================================
\section{Discussion and Limitations}\label{sec:discussion}
%==================================================================

\subsection{What is derived}

\textbf{Exact results:}
\begin{itemize}
    \item The continuity equation for the current velocity $v$ (Proposition~\ref{prop:continuity}).
    \item Irrotationality of $v$ and the existence of the phase function $S$ (Proposition~\ref{prop:irrotational}), a structural consequence of the closure gradient drift.
    \item The closure Madelung pair (Proposition~\ref{prop:madelung}).
    \item The closure-Nelson wave equation~\eqref{eq:closure_schrodinger} (Proposition~\ref{prop:wave}): exact, but nonlinear.
\end{itemize}

\textbf{Exact but regime-specific:}
\begin{itemize}
    \item At exact equilibrium, $U[\Psi_{\mathrm{st}}] = \sigma^2\VW$ and the wave equation is exactly the Schr\"odinger equation with the Witten Hamiltonian (Proposition~\ref{prop:linearisation}, stage~1).
\end{itemize}

\textbf{Leading-order / approximate:}
\begin{itemize}
    \item Near-equilibrium linearisation to $i\partial_t\Psi = H_{\mathrm{W}}\Psi$, valid to leading order in perturbations around $\rho_{\mathrm{st}}$ (Proposition~\ref{prop:linearisation}, stage~2).
    \item The effective parameters $m_{\mathrm{eff}} = 1/\sigma^2$ and $V_{\mathrm{eff}} = \VW$ are read off within the linearised regime (Remark~\ref{rem:parameters}).
\end{itemize}

\subsection{What is conditional}

\begin{itemize}
    \item Exact Schr\"odinger form holds at equilibrium, and the near-equilibrium linearisation is valid to leading order for sufficiently small perturbations. Away from equilibrium, the exact closure-Nelson dynamics remains nonlinear through the correction $\sigma^4\Delta|\Psi|/|\Psi|$.
    \item The linearisation requires the perturbation $\varepsilon\eta$ to be small relative to $\Psi_{\mathrm{st}}$. The range of validity is not quantified.
    \item The identification $\hbar_{\mathrm{eff}} = \sigma^2$ (before unit absorption) is a structural correspondence within this construction, not a derived physical identification.
\end{itemize}

\subsection{What is not achieved}

\begin{itemize}
    \item \textbf{No exact global Schr\"odinger recovery.} The full closure-Nelson equation is nonlinear; the Schr\"odinger equation is linear.
    \item \textbf{No Hilbert-space axioms.} The construction produces a wave equation, not a full quantum-mechanical framework with measurement postulates.
    \item \textbf{No physical derivation of the backward process.} The time-reversed process is defined mathematically; whether it has independent physical content within the closure programme is not addressed.
    \item \textbf{The quantum-potential sign reversal.} The closure Madelung equation has the quantum potential with the opposite sign to standard quantum mechanics. This is absorbed into the nonlinear potential $U$ but represents a genuine structural difference.
    \item \textbf{No exact identification with standard quantum mechanics away from equilibrium.} The paired closure dynamics agrees with linear Schr\"odinger evolution only in the equilibrium and near-equilibrium regime.
\end{itemize}

\subsection{Open questions}

\begin{itemize}
    \item Can the nonlinear correction be absorbed by a redefinition or gauge transformation?
    \item Does the nonlinear Schr\"odinger equation~\eqref{eq:closure_schrodinger} have additional conserved quantities?
    \item Can the range of validity of the linearisation be characterised (e.g.\ in terms of R\'enyi divergence from equilibrium)?
    \item Is there a variational principle for the closure Madelung pair that would select the dynamics independently of the Nelson construction?
\end{itemize}

%==================================================================
\section{Conclusion}\label{sec:conclusion}
%==================================================================

Paper~XVIII changes the location of the remaining problem in the VFD quantum-recovery programme.

\textbf{Exact results.} Starting from the forward closure Langevin process (Paper~XIV) and its time-reversal, the Nelson pairing produces:
\begin{enumerate}
    \item an exact continuity equation for the current velocity,
    \item an automatically irrotational phase function (forced by the closure gradient structure),
    \item an exact closure Madelung pair with a quantum-potential-like term,
    \item a nonlinear Schr\"odinger-type wave equation~\eqref{eq:closure_schrodinger} with the correct imaginary kinetic structure.
\end{enumerate}

\textbf{Equilibrium Schr\"odinger recovery.} At the Paper~XIV equilibrium, the nonlinear potential evaluates exactly to $\sigma^2\VW$, and the wave equation becomes exactly the Schr\"odinger equation with the Witten Hamiltonian:
\begin{equation}
    i\,\partial_t\Psi = H_{\mathrm{W}}\Psi = \left(-\frac{\sigma^2}{2}\Delta + \VW\right)\Psi
\end{equation}
The Witten Hamiltonian --- which appeared in Paper~XVI as the dissipative real part of the drift-free generator --- now reappears as the quantum Hamiltonian governing oscillatory dynamics. For perturbations around equilibrium, this Schr\"odinger equation persists to leading order.

\textbf{What remains.} The closure-Nelson construction bypasses the single-generator obstruction identified in Paper~XVII: the imaginary kinetic term emerges from the forward/backward pairing, and the recovered linear Schr\"odinger dynamics is exact on the equilibrium solution itself and formal (leading-order) near it. The remaining discrepancy is confined to the potential sector: an equilibrium-centred residual term (Section~\ref{sec:residual}) that vanishes at equilibrium and is $O(\varepsilon)$ for near-equilibrium perturbations.

This shifts the problem from the existence of quantum dynamics to the interpretation of this residual term. The closure framework therefore reproduces standard linear quantum mechanics in a well-defined equilibrium-centred regime while yielding a controlled, structurally derived deviation away from that regime. Determining the nature of this deviation --- whether it is a physical nonlinear correction, a gauge artefact, or a structurally meaningful quantity --- is the next stage of the programme.

\textbf{Programme arc.}
\begin{itemize}
    \item Papers~XII--XIV: dissipative closure dynamics,
    \item Paper~XV: phase and interference,
    \item Paper~XVI: evolution equation and structural gap,
    \item Paper~XVII: dissipative obstruction theorem,
    \item Paper~XVIII: Nelson pairing, near-equilibrium Schr\"odinger recovery, and identification of the nonlinear residual as the next open problem.
\end{itemize}

%==================================================================
\bibliographystyle{plainnat}
\bibliography{references}

\end{document}
